\clearpage
\section{Chapter 8: Dark Matter}

\subsection{Questions}

\begin{enumerate}[label=8.\arabic*]
  \item \label{q:8.1} \textbf{Galaxy Rotation Curves}
  
  For a spherical mass distribution $\rho(r)\propto r^{-n}$ with $0<n<3$:
  \begin{enumerate}[label=(\alph*)]
    \item Show that $M(r)\propto r^{3-n}$ and $v(r)\propto r^{(2-n)/2}$.
    \item What value of $n$ gives a flat rotation curve $v\approx \text{const}$?
    \item Explain why the visible matter profile of galaxies typically predicts a declining $v(r)$.
  \end{enumerate}

  \item \label{q:8.2} \textbf{MOND Deep-Regime Scaling}
  
  In MOND, $a\,\mu(a/a_0)=a_N$ and for $a\ll a_0$, $\mu(x)\approx x$.
  \begin{enumerate}[label=(\alph*)]
    \item Derive $v^4 = G M a_0$ for circular orbits.
    \item If $a_0 = 1.2\times10^{-10}\,\mathrm{m/s^2}$ and $M=10^{11}M_\odot$, estimate $v$.
  \end{enumerate}

  \item \label{q:8.3} \textbf{WIMP Freeze-Out and Relic Density}
  
  Use the approximate relic abundance relation
  \[
    \Omega_\chi h^2 \approx 0.1\left(\frac{x_{\mathrm{FO}}}{20}\right)
    \left(\frac{g_*}{80}\right)^{-1/2}
    \left(\frac{\langle\sigma v\rangle}{3\times10^{-26}\,\mathrm{cm^3/s}}\right)^{-1}.
  \]
  \begin{enumerate}[label=(\alph*)]
    \item For $x_{\mathrm{FO}}=25$ and $g_*=90$, find $\langle\sigma v\rangle$ that gives $\Omega_\chi h^2=0.12$.
    \item Comment on how the required $\langle\sigma v\rangle$ compares with the weak scale.
  \end{enumerate}

  \item \label{q:8.4} \textbf{Gamma-Ray Flux and J-Factor}
  
  For annihilation,
  \[
    \Phi_\gamma(E,\psi)=\frac{1}{2}\frac{\langle\sigma v\rangle}{4\pi m_\chi^2}
    \frac{dN_\gamma}{dE}\int_{\text{l.o.s.}}\rho^2(r,\psi)\,dl.
  \]
  \begin{enumerate}[label=(\alph*)]
    \item Identify the ``particle physics'' and ``astrophysical'' factors.
    \item Explain why the flux is proportional to $\rho^2$ rather than $\rho$.
  \end{enumerate}
\end{enumerate}

\bigskip
\bigskip
\bigskip

\subsection{Solutions}

\subsubsection*{Solution 8.1: Galaxy Rotation Curves}

\begin{enumerate}[label=(\alph*)]
  \item
  \[
    M(r)=\int_0^r 4\pi r'^2\rho(r')dr'\propto \int_0^r r'^{2-n}dr'\propto r^{3-n}.
  \]
  Then
  \[
    v^2=\frac{GM(r)}{r}\propto r^{2-n}\Rightarrow v\propto r^{(2-n)/2}.
  \]

  \item Flat rotation curve requires $v\propto r^0$, so $2-n=0$ and $n=2$.

  \item Visible matter falls off faster than $r^{-2}$ in the outer disk, so $n>2$;
  therefore $v(r)$ is predicted to decrease with $r$, unlike the observed flat curves.
\end{enumerate}

\subsubsection*{Solution 8.2: MOND Deep-Regime Scaling}

\begin{enumerate}[label=(\alph*)]
  \item In deep MOND, $a\mu(a/a_0)\approx a(a/a_0)=a_N=GM/r^2$, so
  \[
    a^2=a_0\frac{GM}{r^2}.
  \]
  With $a=v^2/r$:
  \[
    \left(\frac{v^2}{r}\right)^2=a_0\frac{GM}{r^2}\Rightarrow v^4=GMa_0.
  \]

  \item $M=10^{11}M_\odot\approx 1.99\times10^{41}\,\mathrm{kg}$. Then
  \[
    v=(GMa_0)^{1/4}\approx \left((6.67\times10^{-11})(1.99\times10^{41})(1.2\times10^{-10})\right)^{1/4}
    \approx 2\times10^5\,\mathrm{m/s}\approx 200\,\mathrm{km/s}.
  \]
\end{enumerate}

\subsubsection*{Solution 8.3: WIMP Freeze-Out and Relic Density}

\begin{enumerate}[label=(\alph*)]
  \item Rearranging,
  \[
    \langle\sigma v\rangle \approx 3\times10^{-26}\,\frac{x_{\mathrm{FO}}}{20}
    \left(\frac{g_*}{80}\right)^{-1/2}\left(\frac{0.1}{\Omega_\chi h^2}\right)\,\mathrm{cm^3/s}.
  \]
  Plugging in $x_{\mathrm{FO}}=25$, $g_*=90$, $\Omega_\chi h^2=0.12$:
  \[
    \langle\sigma v\rangle \approx 3\times10^{-26}\times 1.25\times\left(\frac{90}{80}\right)^{-1/2}\times\frac{0.1}{0.12}
    \approx 2.9\times10^{-26}\,\mathrm{cm^3/s}.
  \]

  \item The value is close to the weak-scale annihilation cross section, illustrating the
  ``WIMP miracle.''
\end{enumerate}

\subsubsection*{Solution 8.4: Gamma-Ray Flux and J-Factor}

\begin{enumerate}[label=(\alph*)]
  \item Particle-physics factor:
  \[
    \frac{1}{2}\frac{\langle\sigma v\rangle}{4\pi m_\chi^2}\frac{dN_\gamma}{dE}.
  \]
  Astrophysical factor (J-factor):
  \[
    J(\psi)=\int_{\text{l.o.s.}}\rho^2(r,\psi)\,dl.
  \]

  \item Annihilation requires two DM particles, so the rate scales as
  $n_\chi^2\propto \rho^2$.
\end{enumerate}
